
\documentclass{article}
\usepackage{colortbl}
\usepackage{makecell}
\usepackage{multirow}
\usepackage{supertabular}

\begin{document}

\newcounter{utterance}

\centering \large Interaction Transcript for game `wordle\_withclue', experiment `medium\_frequency\_words\_clue\_no\_critic', episode 6 with Qwen3.5{-}4B.
\vspace{24pt}

{ \footnotesize  \setcounter{utterance}{1}
\setlength{\tabcolsep}{0pt}
\begin{supertabular}{c@{$\;$}|p{.15\linewidth}@{}p{.15\linewidth}p{.15\linewidth}p{.15\linewidth}p{.15\linewidth}p{.15\linewidth}}
    \# & \multicolumn{2}{c}{Player} && \multicolumn{2}{c}{Game Master} \\
    \hline

    \theutterance \stepcounter{utterance}  
    & & & \multicolumn{4}{p{0.6\linewidth}}{
        \cellcolor[rgb]{0.9,0.9,0.9}{
            \makecell[{{p{\linewidth}}}]{
                \texttt{\tiny{[P1$\langle$GM]}}
                \texttt{You are a language wizard who likes to guess words by using the given rules.} \\
\\ 
\texttt{Welcome to Wordle! You have six attempts to guess the target word, a valid English word of five lowercase letters (a{-}z). Please use the tags "explanation:" and "guess:" to provide a concise explanation for each guess.} \\
\\ 
\texttt{To help you make an informed guess, you will receive a clue for the word, such as} \\
\texttt{clue: snowy white.} \\
\\ 
\texttt{Here is an example guess based on the clue:} \\
\texttt{explanation: In the fairy tail Snow White, the girl is killed because she eats a poisoned apple. And the word apple has 5 letters.} \\
\texttt{guess: apple} \\
\\ 
\texttt{After each guess, your answer will be validated, and you will receive feedback indicating which letters are correct (green), which letters are correct but in the wrong position (yellow), and which letters are incorrect (red). This feedback can be useful in determining which letters to include or exclude in your next guess.} \\
\\ 
\texttt{For example, the feedback for "apple" might be:} \\
\texttt{guess\_feedback: a<yellow> p<yellow> p<green> l<yellow> e<red>} \\
\\ 
\texttt{The explanation should contain details about how both the guess\_feedback and clue work together in tandem to formulate a new guess.} \\
\\ 
\texttt{Do not generate any other text.} \\
\\ 
\texttt{Let's begin with your first guess.} \\
\\ 
\texttt{clue: like algae} \\
            }
        }
    }
     \\ \\

    \theutterance \stepcounter{utterance}  
    & \multicolumn{4}{p{0.6\linewidth}}{
        \cellcolor[rgb]{0.561,0.737,0.561}{
            \makecell[{{p{\linewidth}}}]{
                \texttt{\tiny{[P1$\rangle$GM]}}
                \texttt{explanation: The clue "like algae" suggests a word related to aquatic plants or greenery. "Seaweed" is a common term for algae, but it has six letters. "Kelp" is a type of algae but only has four letters. "Green" describes the color but isn't a specific type of algae. "Moss" is a plant but not typically classified as algae. "Lichen" is a composite organism, not strictly algae. "Diatom" is a type of algae but has six letters. "Phytopl" is not a word. "Chloro" is a prefix. Let's reconsider common 5{-}letter words. "Seagr" isn't a word. "Algae" itself is 5 letters and matches the clue perfectly, but usually, the clue describes the word, not the word itself. Is there another word? "Marine" is 6 letters. "Ocean" is 5 letters but not algae. "Fresh" is 5 letters. "Liquid" is 6. "Water" is 5. Maybe the word is "seaweed" shortened? No. Let's think of synonyms or related terms. "Green" is 4. "Leaf" is 4. "Root" is 4. "Stem" is 4. "Branch" is 6. "Bloom" is 5. "Flora" is 5. "Flora" relates to plants including algae. "Fauna" is 5. "Zoo" is 3. "Bio" is 3. "Life" is 4. "Organ" is 5. "Cell" is 4. "Nucle" is 5. "Cytopl" is 6. "Mitochondria" is too long. Let's try "algae" again. If the clue is "like algae", maybe the word is "seaweed" and I miscounted? S{-}E{-}A{-}W{-}E{-}E{-}D is 7. "Kelp" is 4. "Sargassum" is 9. "Ulva" is 4. "Gracilaria" is 10. "Porphyra" is 8. "Chlorella" is 9. "Spirogyra" is 9. "Desmids" is 7. "Volvox" is 6. "Euglena" is 7. "Dinoflagellate" is too long. Maybe the word is "green" and the clue is metaphorical? No. What about "chlorophyll"? Too long. "Pigment"? 7. "Color"? 5. "Hue"? 3. "Tint"? 4. "Shade"? 5. "Dye"? 3. "Stain"? 5. "Wash"?} \\
            }
        }
    }
    & & \\ \\

    \theutterance \stepcounter{utterance}  
    & & & \multicolumn{2}{p{0.3\linewidth}}{
        \cellcolor[rgb]{0.9,0.9,0.9}{
            \makecell[{{p{\linewidth}}}]{
                \texttt{\tiny{[GM$|$GM]}}
                \texttt{Error: The response should contain the 'guess:' keyword exactly once.} \\
            }
        }
    }
    & & \\ \\

    \theutterance \stepcounter{utterance}  
    & & & \multicolumn{2}{p{0.3\linewidth}}{
        \cellcolor[rgb]{0.9,0.9,0.9}{
            \makecell[{{p{\linewidth}}}]{
                \texttt{\tiny{[GM$|$GM]}}
                \texttt{game\_result = ABORT} \\
            }
        }
    }
    & & \\ \\

\end{supertabular}
}

\end{document}
